\chapter{Notation and Symbols}
\label{app:notation}

This appendix provides a consolidated reference for all symbols,
conventions, and abbreviations used in the monograph. Symbols are
grouped by category; within each group they appear roughly in the order
of first use.


%% ================================================================
\section{General conventions}
\label{sec:conventions}

\begin{center}
\begin{tabular}{ll}
\toprule
Convention & Choice in this book \\
\midrule
Time-harmonic factor & $e^{-i\omega t}$ \\
Fourier sign & $\tilde{f}(\omega)
 = \int_{-\infty}^{\infty} f(t)\,e^{+i\omega t}\,\dd t$ \\
SI units throughout & \\
Vacuum permittivity & $\varepsilon_0$; relative: $\epsr$ \\
Vacuum permeability & $\mu_0$; relative: $\mur$ \\
Speed of light & $c = 1/\sqrt{\varepsilon_0\mu_0}$ \\
Passive medium & $\Im\epsr(\omega) > 0$ \\
Gain medium & $\Im\epsr(\omega) < 0$ \\
QNM decay & $\omtil = \omega_r - i\gamma$, $\gamma > 0$ \\
Quality factor & $\Qfac = \omega_r / (2\gamma)$ \\
\bottomrule
\end{tabular}
\end{center}


%% ================================================================
\section{Electromagnetic fields and sources}
\label{sec:em-symbols}

\begin{center}
\begin{tabular}{lll}
\toprule
Symbol & Name & SI Unit \\
\midrule
$\EE(\rr, t)$ & Electric field & V\,m$^{-1}$ \\
$\HH(\rr, t)$ & Magnetic field & A\,m$^{-1}$ \\
$\DD(\rr, t)$ & Electric displacement & C\,m$^{-2}$ \\
$\BB(\rr, t)$ & Magnetic flux density & T \\
$\JJ(\rr, t)$ & Current density & A\,m$^{-2}$ \\
$\rho(\rr, t)$ & Charge density & C\,m$^{-3}$ \\
$\SS = \EE \times \HH$ & Poynting vector & W\,m$^{-2}$ \\
$u$ & Electromagnetic energy density & J\,m$^{-3}$ \\
$\sigma$ & Conductivity & S\,m$^{-1}$ \\
\bottomrule
\end{tabular}
\end{center}


%% ================================================================
\section{Material and structural parameters}
\label{sec:material-symbols}

\begin{center}
\begin{tabular}{lll}
\toprule
Symbol & Meaning & Introduced \\
\midrule
$\epsr(\rr,\omega)$ & Complex relative permittivity & Ch.~\ref{ch:light-eigenvalue} \\
$\mur(\rr,\omega)$ & Complex relative permeability & Ch.~\ref{ch:light-eigenvalue} \\
$n = \sqrt{\epsr\mur}$ & Complex refractive index & Ch.~\ref{ch:loss-gain-dispersion} \\
$\omega_0$ & Oscillator resonance frequency & Ch.~\ref{ch:loss-gain-dispersion} \\
$\omega_p$ & Plasma frequency & Ch.~\ref{ch:loss-gain-dispersion} \\
$\gamma$ & Damping rate / gain--loss contrast & Chs.~\ref{ch:loss-gain-dispersion}, \ref{ch:exceptional-points} \\
$F$ & Oscillator strength & Ch.~\ref{ch:loss-gain-dispersion} \\
$a$, $\aa_j$ & Lattice constant, lattice vectors & Ch.~\ref{ch:periodicity} \\
$\GG$ & Reciprocal lattice vector & Ch.~\ref{ch:periodicity} \\
$d$ & Layer thickness & App.~\ref{app:transfer} \\
\bottomrule
\end{tabular}
\end{center}


%% ================================================================
\section{Operators and eigenproblems}
\label{sec:operator-symbols}

\begin{center}
\begin{tabular}{lll}
\toprule
Symbol & Meaning & Introduced \\
\midrule
$\hat{\Theta}$ & Maxwell curl--curl operator & Ch.~\ref{ch:light-eigenvalue} \\
$\omtil_n = \omega_n - i\gamma_n$ & Complex QNM eigenfrequency & Ch.~\ref{ch:qnm} \\
$\tilde{\EE}_n$, $\tilde{\HH}_n$ & QNM field profiles & Ch.~\ref{ch:qnm} \\
$|r_j\rangle$, $\langle l_j|$ & Right, left eigenvectors & App.~\ref{app:nonnormal} \\
$N_j$ & Biorthogonal normalization & Ch.~\ref{ch:qnm}, App.~\ref{app:nonnormal} \\
$K_j$ & Petermann factor & Ch.~\ref{ch:qnm}, App.~\ref{app:nonnormal} \\
$r_j$ & Phase rigidity & App.~\ref{app:nonnormal} \\
$J_m(\lambda)$ & Jordan block of size $m$ & App.~\ref{app:nonnormal} \\
$\sigma_\epsilon(A)$ & $\epsilon$-pseudospectrum & App.~\ref{app:nonnormal} \\
\bottomrule
\end{tabular}
\end{center}


%% ================================================================
\section{Scattering and transfer matrices}
\label{sec:scattering-symbols}

\begin{center}
\begin{tabular}{lll}
\toprule
Symbol & Meaning & Introduced \\
\midrule
$\Smat(\omega)$ & Scattering matrix & Ch.~\ref{ch:boundaries} \\
$\Tmat$ & Transfer matrix & Ch.~\ref{ch:pt-symmetric}, App.~\ref{app:transfer} \\
$t$, $t'$ & Transmission coefficients & App.~\ref{app:transfer} \\
$r_L$, $r_R$ & Reflection coefficients (left, right) & App.~\ref{app:transfer} \\
$\det\Smat$ & Determinant (CPA condition) & Ch.~\ref{ch:poles-zeros} \\
\bottomrule
\end{tabular}
\end{center}


%% ================================================================
\section{Topology and band theory}
\label{sec:topology-symbols}

\begin{center}
\begin{tabular}{lll}
\toprule
Symbol & Meaning & Introduced \\
\midrule
$\kk$, $k$ & Bloch wave vector & Ch.~\ref{ch:periodicity} \\
$\beta = e^{ik a}$ & Bloch phase factor & Ch.~\ref{ch:skin-effect} \\
$\mathcal{A}_n(\kk)$ & Berry connection & Ch.~\ref{ch:berry-maxwell} \\
$\mathcal{F}_{xy}$ & Berry curvature & Ch.~\ref{ch:berry-maxwell} \\
$C$ & Chern number & Ch.~\ref{ch:berry-maxwell} \\
$\wnum$ & Spectral winding number & Ch.~\ref{ch:complex-bands} \\
$q$ & BIC topological charge & Ch.~\ref{ch:bic-topology} \\
$E_{\mathrm{ref}}$ & Point-gap reference energy & Ch.~\ref{ch:point-gap} \\
GBZ & Generalized Brillouin zone & Ch.~\ref{ch:skin-effect} \\
$P^{\mathrm{bi}}$ & Biorthogonal polarization & Ch.~\ref{ch:bbc-rebuilt} \\
\bottomrule
\end{tabular}
\end{center}


%% ================================================================
\section{Abbreviations}
\label{sec:abbreviations}

\begin{center}
\begin{tabular}{ll}
\toprule
Abbreviation & Meaning \\
\midrule
BIC & Bound state in the continuum \\
BBC & Bulk-boundary correspondence \\
BdG & Bogoliubov--de Gennes \\
BZ & Brillouin zone \\
BZF & Brillouin zone folding \\
CPA & Coherent perfect absorption \\
EP & Exceptional point \\
EP$_n$ & Exceptional point of order $n$ \\
FDTD & Finite-difference time-domain \\
FEM & Finite-element method \\
GBZ & Generalized Brillouin zone \\
LDOS & Local density of states \\
NH & Non-Hermitian \\
NHSE & Non-Hermitian skin effect \\
OBC & Open boundary conditions \\
PBC & Periodic boundary conditions \\
PEC & Perfect electric conductor \\
PMC & Perfect magnetic conductor \\
PML & Perfectly matched layer \\
$\PT$ & Parity--time \\
PWE & Plane-wave expansion \\
QNM & Quasinormal mode \\
SNR & Signal-to-noise ratio \\
SSH & Su--Schrieffer--Heeger \\
TCMT & Temporal coupled-mode theory \\
\bottomrule
\end{tabular}
\end{center}


%% ================================================================
\section{Sign conventions and equivalent formulations}
\label{sec:sign-conventions}

Non-Hermitian photonics inherits two competing sign conventions from
the optics and physics communities, and careless mixing of conventions
is a common source of sign errors in the literature. This section
records the choices made in this book and their relationship to the
alternatives.

\subsection{Time-harmonic convention}
\label{ssec:time-convention}

The convention throughout is $e^{-i\omega t}$, following the majority of the
optics and photonics literature. The condensed-matter and high-energy
communities often use $e^{+i\omega t}$. The two conventions are
related by complex conjugation of all frequency-domain quantities:
\begin{center}
\begin{tabular}{lcc}
 \toprule
 \textbf{Quantity} & $e^{-i\omega t}$ (this book) & $e^{+i\omega t}$ \\
 \midrule
 Passive loss & $\Im\epsr > 0$ & $\Im\epsr < 0$ \\
 Gain & $\Im\epsr < 0$ & $\Im\epsr > 0$ \\
 QNM decay & $\omtil = \omega_r - i\gamma$ & $\omtil = \omega_r + i\gamma$ \\
 Outgoing wave & $e^{+ikr}$ & $e^{-ikr}$ \\
 Kramers--Kronig & $\epsr(\omega)$ analytic in upper half-plane & lower half-plane \\
 \bottomrule
\end{tabular}
\end{center}
When importing results from a paper that uses the opposite convention,
replace $\omega \to -\omega^*$ (equivalently, $i \to -i$ in all
material parameters and eigenvalues). The physical observables---power,
transmission, quality factor---are convention-independent.

\subsection{Three formulations of the Petermann factor}
\label{ssec:petermann-conventions}

The Petermann factor quantifies the non-orthogonality of QNMs and
appears in three equivalent forms in the literature. All three are
used in this book at different points; we record the equivalences here
for reference.

\paragraph{Formulation 1: unconjugated-norm ratio
(Chapter~\ref{ch:qnm}).}
\begin{equation}
 K_n = \frac{
 \langle \tilde{\EE}_n^*, \tilde{\EE}_n^* \rangle
 \cdot
 \langle \tilde{\EE}_n, \tilde{\EE}_n \rangle
 }{
 |\langle \tilde{\EE}_n, \tilde{\EE}_n \rangle_{\mathrm{uc}}|^2
 },
 \label{eq:petermann-form1}
\end{equation}
where $\langle\cdot,\cdot\rangle$ is the standard conjugated inner
product and $\langle\cdot,\cdot\rangle_{\mathrm{uc}}$ is the
unconjugated (bilinear) inner product.

\paragraph{Formulation 2: biorthogonal overlap
(Appendix~\ref{app:nonnormal}).}
\begin{equation}
 K_n = \frac{1}{|\langle l_n | r_n \rangle|^2},
 \label{eq:petermann-form2}
\end{equation}
where $|r_n\rangle$ and $\langle l_n|$ are the right and left
eigenvectors first normalized to unit norm in the chosen conjugated
inner product. Equivalently, without that convention one should use
$K_n=\langle r_n|r_n\rangle\langle l_n|l_n\rangle/
|\langle l_n|r_n\rangle|^2$. This formulation makes the
connection to non-normality most transparent: $K_n = 1$ for a normal
operator and $K_n \to \infty$ at an EP where
$\langle l_n | r_n \rangle \to 0$.

\paragraph{Formulation 3: phase rigidity
(Chapter~\ref{ch:exceptional-points}).}
\begin{equation}
 K_n = \frac{1}{|r_n^{\mathrm{ph}}|^2},
 \qquad
 r_n^{\mathrm{ph}} = \frac{\langle \tilde{\EE}_n, \tilde{\EE}_n
 \rangle_{\mathrm{uc}}}
 {\langle \tilde{\EE}_n^*, \tilde{\EE}_n \rangle},
 \label{eq:petermann-form3}
\end{equation}
where $r_n^{\mathrm{ph}}$ is the phase rigidity measuring the ratio of
the unconjugated to the conjugated inner product. The phase rigidity
equals unity for a standing wave (Hermitian mode) and decreases toward
zero as the mode becomes a travelling wave (non-Hermitian, non-normal).

All three formulations give the same numerical value for a given QNM
when the same bilinear form and normalization convention are used. The
Cauchy--Schwarz inequality guarantees $K_n \ge 1$, with equality when
the corresponding left and right vectors are collinear in the chosen
inner product; normal operators satisfy this for all modes.

\subsection{Gain--loss sign in coupled-mode models}
\label{ssec:gain-loss-sign}

In the $\PT$-symmetric coupled-mode models of
Chapter~\ref{ch:pt-symmetric}, the Hamiltonian is written as
\begin{equation}
 H = \begin{pmatrix}
 \omega_0 + i\gamma & \kappa \\
 \kappa & \omega_0 - i\gamma
 \end{pmatrix}.
 \label{eq:pt-hamiltonian-convention}
\end{equation}
With our $e^{-i\omega t}$ convention, $+i\gamma$ corresponds to
\emph{gain} (the mode amplitude grows as $e^{+\gamma t}$) and
$-i\gamma$ corresponds to \emph{loss}. In the opposite convention
($e^{+i\omega t}$), the signs are reversed: $-i\gamma$ is gain and
$+i\gamma$ is loss.

The $\PT$ phase transition occurs at $\gamma = \kappa$ in both
conventions. Below the transition ($\gamma < \kappa$), the
eigenvalues are real and the system is in the exact $\PT$-symmetric
phase. Above the transition ($\gamma > \kappa$), the eigenvalues form
a complex-conjugate pair and the system is in the broken phase. The
physical content is identical; only the intermediate sign bookkeeping
differs.

\subsection{Transfer-matrix ordering}
\label{ssec:transfer-ordering}

Two conventions exist for the transfer matrix
(Appendix~\ref{app:transfer}):
\begin{enumerate}
 \item \emph{Left-to-right} (this book):
 $\bigl(\begin{smallmatrix} A \\ B \end{smallmatrix}\bigr)_{\text{right}}
 = \Tmat\,
 \bigl(\begin{smallmatrix} A \\ B \end{smallmatrix}\bigr)_{\text{left}}$.
 \item \emph{Right-to-left}:
 $\bigl(\begin{smallmatrix} A \\ B \end{smallmatrix}\bigr)_{\text{left}}
 = \Tmat'\,
 \bigl(\begin{smallmatrix} A \\ B \end{smallmatrix}\bigr)_{\text{right}}$.
\end{enumerate}
The two are related by $\Tmat' = \Tmat^{-1}$. For a reciprocal medium
($\det\Tmat = 1$), the inverse is obtained by swapping the diagonal
elements and negating the off-diagonal elements. When comparing
transfer-matrix results across papers, always check which convention is
used before multiplying matrices.
